\documentclass[11pt]{article}
\usepackage[margin=1in]{geometry}
\usepackage{amsmath,amssymb,amsthm,mathrsfs,mathtools,microtype,array,booktabs}
\usepackage[hidelinks]{hyperref}
\usepackage[T1]{fontenc}
\linespread{1.04}

\newtheorem{theorem}{Theorem}[section]
\newtheorem{proposition}[theorem]{Proposition}
\newtheorem{lemma}[theorem]{Lemma}
\newtheorem{corollary}[theorem]{Corollary}
\newtheorem{definition}[theorem]{Definition}
\newtheorem{warning}[theorem]{Warning}
\newtheorem{remark}[theorem]{Remark}
\newtheorem{example}[theorem]{Example}
\newcommand{\A}{\mathbb A}
\newcommand{\Pp}{\mathbb P}
\newcommand{\Spec}{\operatorname{Spec}}
\newcommand{\Aut}{\operatorname{Aut}}
\newcommand{\Isom}{\operatorname{Isom}}
\newcommand{\Hilb}{\operatorname{Hilb}}
\newcommand{\ol}{\overline}

\title{The Distinguished Affine Root at the Boundary:\\
An Adversarial Audit and a Correct Valuative Theorem}
\author{Repository audit draft}
\date{23 July 2026}

\begin{document}
\maketitle

\begin{abstract}
We audit the chain H1--H3 concerning the Hurwitz compactification of
rerooting, stacky extension of a selected zero-fiber root, and descent of
that root to coarse moduli.  We distinguish three assertions that the
current repository conflates: the elementary valuative extension of a
rational point through a finite cover; the quotient-stack description of
choosing one element of an unordered divisor; and the comparison of that
quotient stack with the repository-specific admissible-cover contraction
and regular-reconstruction component.  The first two assertions are proved
here, including arbitrary simultaneous collisions and stabilizers.  The
third is isolated as an explicit comparison hypothesis.  Under that
hypothesis we give a complete proof that the distinguished affine-root
marking extends uniquely over every DVR, as a scheme-valued marking and not
merely as a geometric point.  We then show why the present H1--H3 text does
not prove the comparison hypothesis.  In particular, properness of an
ambient admissible-cover stack, equality of generic function algebras, and
isomorphisms of completed collision charts do not by themselves construct
the required global finite marked cover.  The resulting verdict is that the
valuative statement is correct conditionally and locally well supported,
but H2 and H3 are not presently available as foundational infrastructure.
\end{abstract}

\begin{quote}
\textbf{Post-audit resolution (23 July 2026).}
The companion note
\href{global-comparison-obstruction.tex}{The Missing H1/H2 Comparison
Package} gives a degree-five obstruction to the first comparison morphism.
At the simultaneous boundary $a\to0$, $b\to1$, the admissible target retains
the cross-ratio $x^3/y^2$ that $\overline M_{0,5}$ forgets.  The correct
global object is therefore the normalized graph (locally the normalized
blowup of $(x^3,y^2)$), followed by the normalization/Stein contraction to
the finite incidence.  The valuative and finite-root results proved below
are unchanged.
\end{quote}

\begin{quote}
\textbf{Post-repair resolution (23 July 2026).}
The later
\href{../../extended-geometry/BRANCH_GRAPH_WONDERFUL_PULLBACK.md}
{wonderful-pullback theorem} constructs the normal corrected target graph in
all degrees.  The
\href{../../extended-geometry/SOURCE_GRAPH_FINITE_NORMALIZATION.md}
{source-graph finite-normalization theorem} then pulls back the finite fully
marked admissible-cover branch morphism and normalizes the closure of the
generic polynomial cover.  Its coarse map to the wonderful graph is finite
birational and hence an isomorphism.  This supplies the comparison object
missing at the time of this audit and proves H1-COARSE.  Consequently
corrected H2 and coarse H3 are now unconditional over that graph.  The later
\href{../../extended-geometry/RECURSIVE_RESONANCE_ATLAS.md}
{fs logarithmic H1-STACK theorem} identifies the selected factor with
\((P\oplus_{\mathbb N^E}Q)^{\rm sat}\) and proves that normalization
commutes with the generic-section selection.  Root-residue flag formulas,
contraction maps, and simultaneous inertia characters are consequences of
that log chart.  The conditional
language below is retained as the historical audit of the false root-stable
formulation.
\end{quote}

\tableofcontents

\section{Verdict and logical shape}

Fix a field $k$ of characteristic zero and an integer $N\geq 4$, and put
$n=N-2$.  On the clean seed locus a normalized polynomial
\[
 H(W)=-\frac{W^2(W-1)\prod_{i=1}^{n-1}(W-\rho_i)}
 {\prod_{i=1}^{n-1}(1-\rho_i)}
\]
has an exact double zero at $0$, a distinguished simple zero at $1$, and
$n-1$ further simple nonzero zeros.  The open rerooting cover forgets which
of the $n$ simple zeros was normalized to $1$.  The repository asks whether
the root selected by the regular-reconstruction chart continues uniquely
when some or all of these roots collide.

There is a short correct answer once the right finite morphism has actually
been constructed.  If $R$ is a DVR with fraction field $K$, if
$Y\to X$ is finite, and if an $R$-point of $X$ has a $K$-rational lift to
$Y$, that lift extends uniquely to $R$ whenever $R$ is integrally closed.
This is the entire scheme-valued valuative argument.  It works for a
nonreduced collision fiber and does not infer uniqueness from the number of
geometric points.

The hard point is therefore upstream:
\begin{quote}
construct a global finite marked object over the precise compactification,
and identify its generic section with the distinguished
regular-reconstruction component.
\end{quote}
The current H2/H3 proof assumes this in the phrases ``the contracted ambient
root cover'' and ``the closure of the distinguished component.''  H1 cites
proper admissible-cover technology but does not define the repository
subfunctor, construct its contraction, or prove that it is the quotient
compactification used in the later invariant-ring calculation.

Our conclusion is consequently two-part.
\begin{enumerate}
\item The quotient-stack and coarse collision model has the desired DVR
extension theorem.  We prove it in Sections~\ref{sec:functors}--\ref{sec:dvr}.
\item Application to the claimed Hurwitz--LL compactification is conditional
on the global comparison package in Definition~\ref{def:comparison}.  The
repository does not yet prove that package; Section~\ref{sec:audit} records
the precise gaps.
\end{enumerate}

This distinction matters for OP-LR-NE.  Unique specialization of one point
of a finite cover is not properness of a relative $\Isom$ functor.  Even
after the comparison package is proved, OP-LR-NE still requires a bounded,
representable, closed extension of isomorphisms and a proof that its
projection is proper and quasi-finite.

\section{The moduli functors}\label{sec:functors}

We define all objects used in the valuative statement.  The definitions are
deliberately independent of any chosen presentation by coefficients.

\subsection{Stable root configurations}

Let $L=\{0,\infty,1,\ldots,n\}$, a set of $n+2=N$ labels.  Write
$\ol M_{0,L}$ for the fine moduli scheme of stable genus-zero curves with
pairwise disjoint sections $(p_\ell)_{\ell\in L}$.  Thus an object over a
scheme $T$ is a proper flat nodal curve $C\to T$ of arithmetic genus zero,
together with sections through the smooth locus, such that every geometric
component has at least three special points.  The symmetric group $S_n$
permutes the labels $1,\ldots,n$ and fixes $0,\infty$.

For a finite constant group $G$ acting on a scheme $U$, we use the following
functorial definition of the quotient stack.  An object of $[U/G](T)$ is a
right $G$-torsor $P\to T$ in the etale topology together with a
$G$-equivariant morphism $P\to U$.  A morphism is an isomorphism of torsors
compatible with the maps to $U$.  Set
\[
 \mathscr B_n=[\ol M_{0,L}/S_n],\qquad
 \mathscr S_n=[\ol M_{0,L}/S_{n-1}],                 \tag{2.1}
\]
where $S_{n-1}\subset S_n$ fixes the label $1$.  We call $\mathscr B_n$ the
unmarked collision-separating root stack and $\mathscr S_n$ the marked one.
The morphism
\[
 q:\mathscr S_n\longrightarrow\mathscr B_n           \tag{2.2}
\]
forgets which simple-root label is $1$.

The word ``marked'' in (2.1) does not mean that every label is globally
ordered on the base.  It means that one element of the unordered degree-$n$
divisor is selected; the other $n-1$ elements remain unordered.  Stable
bubbles keep all labelled points disjoint even when their contraction has a
multiple root.

\subsection{The open normalized-seed functors}

Let $\mathscr U_n\subset\mathscr B_n$ be the open on which the source curve
is smooth and the points $p_0,p_\infty,p_1,\ldots,p_n$ are distinct.  On an
etale cover where the $n$ simple points are labelled, the unique coordinate
$W$ satisfying $p_0=0$ and $p_\infty=\infty$ is still subject to scaling.
Choosing one simple point as $p_1=1$ removes that scaling.  The corresponding
polynomial is, up to the endpoint derivative normalization,
\[
 W^2\prod_{i=1}^{n}(W-p_i).
\]
We denote the marked open by $\mathscr U_n^{\rm mark}=q^{-1}(\mathscr U_n)$.
The repository's clean normalized seed space maps to
$\mathscr U_n^{\rm mark}$ by sending $H$ to its zero divisor with $0$,
$\infty$, and the root $1$ distinguished.

This map is elementary on the open.  Nothing in its formula constructs an
extension to a boundary compactification: when points collide, a stable
configuration acquires bubbles, while a coefficient model contracts them.
Relating those two operations is a separate theorem.

\subsection{Admissible polynomial covers}

We next give a sufficiently precise functor for the external compactification
technology.  Fix the ramification profile consisting of total ramification
over a target section $b_\infty$ and profile $(2,1^n)$ over a target section
$b_0$.  Let $\mathscr H_N^{\rm adm}(T)$ be the groupoid of data
\[
 (C\xrightarrow{f}P; b_0,b_\infty; p_0,p_\infty; D),              \tag{2.3}
\]
where $P\to T$ is a stable pointed nodal genus-zero target (possibly twisted
at nodes and markings), $C\to T$ is a nodal source, $f$ is a finite flat
degree-$N$ admissible cover, $p_\infty$ is the totally ramified point over
$b_\infty$, $p_0$ has ramification index two over $b_0$, and $D$ is the
residual degree-$n$ unramified divisor in $f^{-1}(b_0)$.  At nodes the two
branch indices agree, and the usual stability condition makes automorphism
groups finite.  Branch data away from $b_0,b_\infty$ must also be fixed in a
given connected Hurwitz component; omitting it does not define a finite-type
Hurwitz stack.

The marked functor $\mathscr H_N^{\rm adm,mark}$ adds a section
$p_1:T\to D$.  In the collision-separating model $D$ is represented by
disjoint markings on the expanded source.  Forgetting $p_1$ gives
\[
 q_{\rm adm}:\mathscr H_N^{\rm adm,mark}
       \longrightarrow\mathscr H_N^{\rm adm}.                         \tag{2.4}
\]
To obtain a compactification of the seed locus one must specify: the Hurwitz
component, all branch markings or their allowed collisions, the target
stability weights, and the closure of the normalized pencil locus.  ACV and
Deopurkar provide ambient stacks after these choices; their properness does
not choose the repository-specific component.

\subsection{The contracted root incidence}

For a monic degree-$n$ polynomial
\[
 F(T)=T^n-e_1T^{n-1}+\cdots+(-1)^ne_n,
\]
the universal root incidence is
\[
 \mathcal R_n=\Spec
 k[e_1,\ldots,e_n,T]/(F(T))
 \longrightarrow \A^n_{e_1,\ldots,e_n}.                            \tag{2.5}
\]
It is finite flat of degree $n$.  It is not etale along the discriminant.
For a family of zero divisors, the contracted root cover is obtained by
base change from (2.5), with the fixed double root removed before forming the
degree-$n$ residual polynomial.  A marking is a section of this finite cover,
equivalently a factorization $F=(T-a)G$ over the base.

The inverse-pencil incidence
\[
 X_H=\Spec k[s,t,W]/(H(W)-sW+t)\longrightarrow\A^2_{s,t}             \tag{2.6}
\]
is instead finite flat of degree $N$.  It parametrizes every point in the
fiber of $H(W)-sW$ over $-t$.  It must not be identified without argument
with the degree-$n$ cover (2.5), which only selects residual simple roots of
the zero fiber after removing the fixed double root.

\section{The quotient-stack theorem}

\begin{proposition}\label{prop:quotient}
The morphism $q$ in (2.2) is representable, finite, and etale of constant
degree $n$.  For every geometric point $x:\Spec\Omega\to\mathscr B_n$, its
base change is the discrete $n$-element scheme of choices of one label.
\end{proposition}

\begin{proof}
Let $X=\ol M_{0,L}$, $G=S_n$, and $H=S_{n-1}$.  Given a scheme
$T\to[X/G]$, let $P\to T$ be the associated $G$-torsor.  A reduction of its
structure group to $H$ is a section of the associated bundle
\[
 P\times^G(G/H)\longrightarrow T.                                    \tag{3.1}
\]
The fiber product $T\times_{[X/G]}[X/H]$ is canonically this associated
bundle, with the compatibility map to $X$ carried along.  Since $G/H$ is a
constant etale scheme with $n$ elements, (3.1) is finite etale of degree
$n$.  The base change by every scheme is a scheme, proving representability.
\end{proof}

\begin{remark}
Proposition~\ref{prop:quotient} is not contradicted by ramification on coarse
spaces.  The stack map remembers stabilizers.  The coarse map forgets them
and is ramified at collisions.  Conflating these two maps is one of the main
hazards in H2/H3.
\end{remark}

\begin{proposition}[relative automorphisms]\label{prop:aut}
Let $y$ be an object of $\mathscr S_n$ with image $x$ in $\mathscr B_n$.
Then
\[
 \Aut_{\mathscr S_n}(y)\longrightarrow\Aut_{\mathscr B_n}(x)
\]
is injective.  Hence an automorphism of a marked lift inducing the identity
on the unmarked object is the identity.
\end{proposition}

\begin{proof}
This is the inertia criterion for representability, or directly the torsor
description: an automorphism of an $H$-reduction that is the identity on the
underlying $G$-torsor is the identity.  Equivalently, the relative inertia of
$q$ is trivial.
\end{proof}

At a collision, the unmarked object can still have nontrivial absolute
stabilizer.  For example, a symmetric cluster of $\mu$ labels can carry an
$S_\mu$ stabilizer, while a choice of one label reduces it to
$S_{\mu-1}$.  Proposition~\ref{prop:aut} says that this does not create a
hidden automorphism of the lift over the fixed unmarked object.  It does not
say that the coarse space remembers the labelled point.

\section{Coarse root charts and simultaneous collisions}

We now prove the invariant-ring calculation used by H3 and state exactly
what it does and does not imply.

\begin{proposition}[one cluster]\label{prop:invariants}
Let $A=k[x_1,\ldots,x_\mu]$, with $S_\mu$ permuting the variables and
$S_{\mu-1}$ fixing $x_1$.  If $e_j$ denotes the elementary symmetric
function in all variables, then
\[
 A^{S_{\mu-1}}
 \cong A^{S_\mu}[T]/
 (T^\mu-e_1T^{\mu-1}+\cdots+(-1)^\mu e_\mu),          \tag{4.1}
\]
with $T\mapsto x_1$.  In particular the coarse forgetful map is finite flat
of degree $\mu$, is generically etale, and has total-collision fiber
$k[T]/(T^\mu)$.
\end{proposition}

\begin{proof}
Put $e'_j=e_j(x_2,\ldots,x_\mu)$.  Then
$e_j=e'_j+Te'_{j-1}$, with $e'_0=1$.  Recursively, every $e'_j$ belongs to
$k[e_1,\ldots,e_\mu,T]$.  Hence $A^{S_{\mu-1}}=k[T,e'_1,\ldots,e'_{\mu-1}]$
is generated over $A^{S_\mu}=k[e_1,\ldots,e_\mu]$ by $T$.  Conversely the
formulas $e_j=e'_j+Te'_{j-1}$ define the inverse change of generators.  The
last recursion equation $e_\mu=Te'_{\mu-1}$ becomes exactly the displayed
monic equation after eliminating $e'_1,\ldots,e'_{\mu-1}$.  Thus the two
presentations are inverse isomorphisms.  Since the equation is monic, the
map is finite flat of rank $\mu$.  Setting every $e_j$ to zero gives the
stated fiber.
\end{proof}

\begin{warning}
The special fiber in Proposition~\ref{prop:invariants} has one geometric
point but $\mu$ scheme-theoretic length.  The fact that its topological space
has one point does not prove uniqueness of a section over a DVR.  Uniqueness
is instead a consequence of separatedness after the generic section has
been fixed; existence uses integrality, as in Section~\ref{sec:dvr}.
\end{warning}

\begin{proposition}[simultaneous multicluster chart]\label{prop:multi}
Suppose the special divisor has disjoint clusters of sizes
$\mu_1,\ldots,\mu_r$, together with simple spectator roots.  After an etale
localization separating the cluster centers, the completed coarse root cover
is a disjoint union, over the possible cluster containing the selected root,
of completed tensor products of the corresponding rings (4.1) with formal
power-series rings in spectator and center parameters.  Its fiber over the
deepest collision stratum has one geometric selected point for each distinct
cluster, and length $\mu_i$ on the component selecting cluster $i$.
\end{proposition}

\begin{proof}
Choose pairwise disjoint etale neighborhoods of the distinct cluster centers.
Hensel factorization writes the residual monic polynomial uniquely as a
product $F_1\cdots F_rF_{\rm sp}$, where $F_i$ has degree $\mu_i$ and
$F_{\rm sp}$ is etale.  The functor of choosing a root is the disjoint union
of the functors of choosing a root of one factor.  Completion commutes with
finite products and with the finite base changes in (4.1).  The parameters
of the other factors are formally independent, giving the asserted completed
tensor products.  Specializing each $F_i$ to $T^{\mu_i}$ gives the length and
geometric-point statements.
\end{proof}

The proposition corrects an imprecision in the phrase ``simultaneous
clusters have one geometric coarse mark.''  There is one point within a
chosen cluster, but different cluster centers remain different geometric
points.  A generic distinguished root determines which component is chosen;
the collision alone does not.

\section{The finite-cover DVR lemma}\label{sec:dvr}

We now prove the scheme-valued extension result in the generality actually
needed.

\begin{lemma}[finite valuative extension]\label{lem:finite-dvr}
Let $R$ be a DVR with fraction field $K$, let $X$ be an algebraic space, and
let $f:Y\to X$ be finite and representable.  Given a commutative generic
diagram
\[
\begin{array}{ccc}
 \Spec K&\longrightarrow&Y\\
 \downarrow&&\downarrow f\\
 \Spec R&\longrightarrow&X,
\end{array}                                                          \tag{5.1}
\]
there is a unique dotted arrow $\Spec R\to Y$ completing the diagram.
The same assertion holds after arbitrary extension of DVRs.
\end{lemma}

\begin{proof}
Base change by $\Spec R\to X$.  Because $f$ is representable and finite, the
fiber product is a finite algebraic space over $\Spec R$, hence an affine
scheme $\Spec B$ with $B$ finite over $R$.  The generic lift is an
$R$-algebra map $B\to K$.  Every element of $B$ is integral over $R$, so its
image in $K$ is integral over $R$.  A DVR is integrally closed in $K$;
therefore the image lies in $R$.  This gives an $R$-algebra map $B\to R$ and
hence the extension.  Since a finite morphism is separated, two extensions
agree if they agree on the dense generic point.  The proof applies verbatim
to any DVR extension.
\end{proof}

\begin{remark}
Etaleness is unnecessary in Lemma~\ref{lem:finite-dvr}.  This is important
for the coarse map, which is ramified at collisions.  Properness alone would
give existence only in the usual valuative sense and may require attention
to field extensions for stacks.  Finiteness plus a $K$-rational lift gives
the stronger extension over the original DVR.
\end{remark}

\begin{corollary}[stacky marking]\label{cor:stack-dvr}
For every DVR map $\Spec R\to\mathscr B_n$, every lift of its generic point
to $\mathscr S_n$ extends uniquely, up to a unique 2-isomorphism, to
$\Spec R\to\mathscr S_n$.  There is no relative stabilizer ambiguity.
\end{corollary}

\begin{proof}
Apply Proposition~\ref{prop:quotient} and
Lemma~\ref{lem:finite-dvr}.  Representability turns the lifting problem into
one for a finite scheme.  Proposition~\ref{prop:aut} makes the relative
2-automorphism group trivial.
\end{proof}

\begin{corollary}[coarse marking]\label{cor:coarse-dvr}
Let $B_n$ and $S_n$ be the coarse spaces of $\mathscr B_n$ and
$\mathscr S_n$, and suppose the induced map $S_n\to B_n$ is finite.  Then a
$K$-rational generic marked lift of an $R$-point of $B_n$ extends uniquely to
an $R$-point of $S_n$.  In the local collision charts this extension is the
unique integral root of the universal monic polynomial that equals the given
generic root.
\end{corollary}

\begin{proof}
Apply Lemma~\ref{lem:finite-dvr} to $S_n\to B_n$.  In chart (4.1), the
generic marking is a map sending $T$ to an element $a\in K$ integral over
$R$.  Thus $a\in R$, and the factorization by $T-a$ extends.  This proves
existence as a scheme-valued section.  Separatedness proves uniqueness.
\end{proof}

\begin{example}
The family $F(T)=T^2-\pi^2$ over a DVR has two generic roots $\pm\pi$ and
both specialize to the unique geometric point of $k[T]/(T^2)$.  Once the
generic root $+\pi$ is fixed, only the section $T\mapsto\pi$ extends it.  This
shows simultaneously why one geometric special point is insufficient for
scheme-valued uniqueness and why the finite-cover lemma gives the right
answer.
\end{example}

\section{The comparison package actually required}

The preceding results apply to the repository only after the following data
are constructed.  We state them as a package so that no conclusion is hidden
inside notation.

\begin{definition}[global comparison package]\label{def:comparison}
A Hurwitz compactification of the normalized seed locus satisfies
\textup{(CMP)} if it comes with:
\begin{enumerate}
\item a proper normal Deligne--Mumford stack $\ol{\mathscr B}^{\rm seed}_N$
containing the unmarked clean seed stack, and a proper normal marked stack
$\ol{\mathscr S}^{\rm seed}_N$ containing the distinguished marked seed
stack;
\item a representable finite morphism
$q_{\rm seed}:\ol{\mathscr S}^{\rm seed}_N\to
\ol{\mathscr B}^{\rm seed}_N$ restricting to the degree-$n$ rerooting cover;
\item a global contraction to a finite degree-$n$ residual-root cover
$\mathcal R^{\rm seed}_N$ and a morphism
$\ol{\mathscr S}^{\rm seed}_N\to\mathcal R^{\rm seed}_N$ whose open
restriction is the selected nonzero zero-fiber root;
\item an identification of the open section in item 3 with the unique
regular-reconstruction component supplied by F2;
\item compatibility of items 1--4 with coarse spaces and base change,
including a finite coarse marked map and the invariant charts of
Propositions~\ref{prop:invariants} and~\ref{prop:multi};
\item if normalized Stein factors or conductor squares are to be used later,
a constructed global comparison morphism identifying those finite models,
not only their generic fields or completed local rings.
\end{enumerate}
\end{definition}

\begin{theorem}[distinguished affine-root extension, conditional form]
\label{thm:main}
Assume \textup{(CMP)}.  Let $R$ be any DVR with fraction field $K$, and let
$\Spec R\to\ol{\mathscr B}^{\rm seed}_N$ be a degeneration whose generic
point lies in the clean locus.  Suppose the distinguished affine root gives
a $K$-lift to $\ol{\mathscr S}^{\rm seed}_N$.  Then:
\begin{enumerate}
\item the lift extends over $R$;
\item it is unique as a scheme-valued lift, up to a unique 2-isomorphism;
\item its image on coarse moduli is the unique extension of the coarse
generic mark;
\item the assertion commutes with arbitrary DVR base change;
\item at simultaneous clusters of sizes $\mu_1,\ldots,\mu_r$, the extension
lies in the component determined by the generic root and has special
scheme-theoretic length $\mu_i$ if that root enters cluster $i$;
\item no automorphism acting trivially on the unmarked degeneration can
change the extended mark.
\end{enumerate}
\end{theorem}

\begin{proof}
Item 1 follows from Lemma~\ref{lem:finite-dvr} applied to
$q_{\rm seed}$.  Item 2 follows from separatedness and representability; the
latter kills relative inertia exactly as in Proposition~\ref{prop:aut}.
By item 5 of (CMP), the induced coarse marked map is finite, so
Corollary~\ref{cor:coarse-dvr} gives item 3.  Both the integral-closure proof
and the uniqueness proof survive DVR extension, giving item 4.

For item 5, pass etale-locally to the Hensel-separated factorization used in
Proposition~\ref{prop:multi}.  The generic root selects exactly one factor.
Its closure cannot jump to another open-and-closed factor, and the chosen
factor has collision algebra $k[T]/(T^{\mu_i})$.  This proves compatibility
with all simultaneous clusters, including mixed pair, triple, and higher
collisions.  Finally, item 6 is the injectivity on stabilizers implied by
representability.  Item 4 of (CMP) identifies this abstractly extended mark
with the distinguished regular-reconstruction component rather than an
arbitrary root label.
\end{proof}

\begin{corollary}
Under \textup{(CMP)}, the closure of the graph of the distinguished section
inside the finite marked cover is isomorphic to the normal DVR base along
every trait.  If the same graph is globally finite and birational over a
normal integral marked coarse compactification, then it is globally
isomorphic to that compactification.
\end{corollary}

\begin{proof}
The trait statement is Theorem~\ref{thm:main}.  The global statement is the
finite-birational criterion.  Notice that global finiteness and the existence
of the morphism are hypotheses; they do not follow from a completed-local
description alone.
\end{proof}

\section{Adversarial audit of H1--H3}\label{sec:audit}

We compare the current manuscript with the hypotheses used above.  The point
is not that the desired comparison is false.  It is that several nontrivial
morphisms are invoked before they are defined.

\subsection{H1: the compactification is not specified enough}

The current proof says that taking the normalization of the closure in
``the proper admissible-cover stack'' gives the desired compactification.
There is no single such stack.  One must fix the Hurwitz component, all
ramification profiles, branch markings, allowed branch collisions, target
weights, and rigidifications.  Deopurkar's construction explicitly varies
with the compactification of the branch divisor.  Ambient properness proves
properness of a closed substack after it has been defined; it does not prove
that the closure equals $[\ol M_{0,N}/S_n]$ or that a coefficient contraction
exists.

The displayed normalized polynomial is determined by its zero divisor on
the smooth locus.  At the boundary, however, a stable pointed source and an
admissible map carry node indices, target expansions, and possibly twisted
structure not encoded by the stabilized zero divisor alone.  A proof of H1
must construct mutually inverse functors, or at least a finite birational
morphism with a verified normal target, between the relevant closure and the
claimed quotient stack.  The manuscript currently passes from one to the
other with the word ``equivalently.''

The LL assertion introduces an additional mismatch.  A polynomial cover is
compactified by branch data, while $\ol M_{0,N}$ in H1 records zero-fiber
source points.  These are related by the polynomial, but their stable limits
need not be identical compactifications.  The stable branch-divisor theorem
extends the branch divisor of a family; it does not identify the closure of
the zero-root configuration with a prescribed Hurwitz component.

\subsection{H2: tautological persistence versus comparison}

On a marked admissible-cover stack, the marked section persists by
definition.  This proves no descent statement and no identification with the
repository reconstruction cover.  The quotient calculation in
Proposition~\ref{prop:quotient} is correct if the marked and unmarked stacks
really are the two quotient stacks in (2.1).  That identification is the
missing H1 comparison, not a consequence of subgroup index.

The formal-comparison proof uses two finite normal algebras with the same
generic function algebra and invokes uniqueness of integral closure.  This
argument is valid only after proving all of the following: the proper
incidence has a coherent Stein algebra finite over the stated base; its
generic algebra is the full algebra of (2.6), with the correct product of
components; the normalization is taken in that algebra; and a global
comparison map to the repository incidence has been constructed.  None of
these is supplied by the phrase ``pull back and mark a point over $-t$.''

Moreover, that phrase defines the degree-$N$ incidence (2.6), whereas H2's
rerooting morphism has degree $n=N-2$.  Removing the fixed length-two zero
cluster and restricting to the zero fiber are operations that must be made
functorial in a degeneration.  Equality of generic fields cannot silently
change the degree.

The completed chart $k[[s,u]]$ is a correct chart for the smooth hypersurface
$X_H$.  An isomorphism of this chart with a contracted admissible chart is
asserted, not derived from an explicit contraction.  Even isomorphisms at
every completion would not automatically glue to a global isomorphism
without a global morphism and suitable finiteness.  The conductor calculation
similarly compares two conductor squares only after the downstairs finite
subalgebras have been identified; equality of normalizations and numerical
conductor exponents is insufficient.

\subsection{H3: the finite-birational argument starts after the gap}

The finite-birational criterion used by H3 is correct.  Its application is
circular.  The proof places a closure $Z$ in ``the contracted finite root
cover,'' declares it finite over the marked coarse compactification, and
then uses normality.  A closure in a finite ambient space is finite only when
that ambient space and its map over the same base have already been
constructed.  This is item 3 of (CMP), the main missing comparison.

There is also a base mismatch in the statement.  The universal root
incidence $S_n\to B_n$ is the coarse marked-to-unmarked map.  A graph of the
tautological selected root over $S_n$ lives in the pullback of the root cover
to $S_n$.  Saying that its closure is finite birational over $S_n$ is useful
only after constructing this graph and proving that the generic
regular-reconstruction component maps to it.  The relative Hilbert scheme
of length-one subschemes of a finite scheme is indeed the finite scheme
itself, but this identity does not choose the desired component or prove the
comparison with the admissible contraction.

Finally, the categorical property of a coarse moduli space only factors
maps from the stack to algebraic spaces.  It does not turn a stack section
of a universal curve into a section of a coarse universal curve when no such
universal curve exists, nor does it prove independence of distinct stacky
lifts of the same coarse trait.  Here independence can be proved from the
explicit finite coarse root incidence and Lemma~\ref{lem:finite-dvr}, but the
coarse map must first be identified globally.

\subsection{Risk matrix}

\begin{center}
\small
\begin{tabular}{>{\raggedright\arraybackslash}p{0.19\textwidth}
                >{\raggedright\arraybackslash}p{0.31\textwidth}
                >{\raggedright\arraybackslash}p{0.39\textwidth}}
\toprule
Claim & What is valid & What is not yet proved \\
\midrule
Marked/unmarked quotient & $[X/S_{n-1}]\to[X/S_n]$ is representable finite
etale of degree $n$ & The seed admissible-cover closures are these quotient
stacks \\
Coarse collision chart & The invariant ring is the universal monic-root
incidence and is ramified & This chart is the coarse contraction of the
repository's global Hurwitz closure \\
Stein comparison & A finite normal model is determined by a fixed generic
algebra & The proper incidence has that exact generic algebra and admits the
claimed global comparison map \\
Completed charts & Hensel factorization separates distinct clusters & The
local identifications arise from and glue to one global contraction \\
Conductor & Formulae can be checked once the finite downstairs algebra is
fixed & Equal exponents alone identify the conductor square \\
DVR uniqueness & A fixed generic lift through a finite map extends uniquely
over a DVR & One geometric point implies a unique scheme-valued lift, or an
ambient proper stack makes a specific $\Isom$ relation proper \\
Stabilizers & Representability kills relative automorphisms & Coarse descent
by itself remembers all stacky choices \\
\bottomrule
\end{tabular}
\end{center}

\section{Consequences for OP-LR-NE}

Theorem~\ref{thm:main}, once unconditional, supplies exactly one ingredient
for OP-LR-NE: a chosen affine sheet cannot jump between components of the
finite marked root cover along a DVR.  It does not establish valuative
properness of an isomorphism relation.

Let $\mathcal X$ and $\mathcal Y$ be two families of marked normalized
incidence covers over a base $T$.  A candidate functor
\[
 \mathcal I=\Isom_T^{\rm dec}(\mathcal X,\mathcal Y)                 \tag{8.1}
\]
must specify which decorations are preserved: the finite cover, Hessian
divisor, affine sheet, conductor, node pairing, and any target/source
polarizations needed for boundedness.  To infer that the closure of
$\mathcal I$ is finite and proper, one still needs:
\begin{enumerate}
\item representability of (8.1), including control of automorphisms at
stacky boundary points;
\item finite type or a boundedness theorem for the allowed isomorphisms;
\item separatedness, usually from a rigid enough decoration;
\item existence of extensions over DVRs, not only extension of the selected
root;
\item quasi-finiteness at every boundary point, not merely generic deck
rigidity.
\end{enumerate}
Generic triviality of a deck group gives generic quasi-finiteness.  It does
not exclude new automorphisms on special fibers.  Properness of the ambient
Hurwitz stack does not imply properness of an open or locally closed
$\Isom$ subfunctor.  Thus OP-LR-NE remains genuinely open after the marking
theorem, exactly as the repository's own final paragraph suggests.

\section{A proof checklist that would close H1--H3}

The following finite list turns the comparison package into verifiable
lemmas.
\begin{enumerate}
\item Define the fully labelled seed Hurwitz functor, including every branch
profile, target marking, stability weight, and rigidification.
\item Construct the morphism from the clean labelled zero-root
configuration to that functor and prove it is an open immersion into one
specified connected component.
\item Extend the construction over $\ol M_{0,N}$, or construct a proper
birational morphism from the normalization of the Hurwitz closure to
$\ol M_{0,N}$ and prove it is an isomorphism.  Check node indices and twisted
stabilizers on every boundary graph.
\item Prove equivariance for $S_n$ and $S_{n-1}$ before taking quotients.
This yields the finite representable etale stack map without ambiguity.
\item Construct, over the unmarked closure, the residual zero-fiber divisor
of degree $n$ after removing the fixed ramification-length-two summand.
Prove it is finite flat and commutes with base change.
\item Construct the contraction from the expanded source to that residual
root cover.  Verify that the selected section maps to its tautological root.
\item Identify the F2 regular-reconstruction component with this section on
the clean open, then take its scheme-theoretic closure.
\item Prove the induced coarse map is finite and establish the local
invariant presentation (4.1) from the global map.  Do not reverse this
implication.
\item For the degree-$N$ inverse incidence and critical incidence, construct
separate proper objects and global Stein comparison morphisms.  Verify the
generic algebras component by component before invoking integral closure.
\item Derive the conductor-square comparison from equality of the downstairs
finite algebras; use completed exponent calculations as checks, not as the
definition of the global isomorphism.
\end{enumerate}

Once items 1--8 are complete, Theorem~\ref{thm:main} gives the requested DVR
extension immediately and without further admissible-cover technology.
Items 9--10 are needed only for the stronger normalized-Stein and conductor
claims bundled into H2.

\section{Final theorem status}

The requested sentence
\begin{quote}
``The distinguished affine-root marking extends uniquely across every DVR
degeneration in the claimed compactification''
\end{quote}
has two readings.  For the explicitly defined quotient compactification
$\mathscr S_n\to\mathscr B_n$, it is proved by
Corollary~\ref{cor:stack-dvr}, Corollary~\ref{cor:coarse-dvr}, and
Proposition~\ref{prop:multi}.  For the repository-specific
Hurwitz/admissible-cover compactification and its regular-reconstruction
contraction, it follows from Theorem~\ref{thm:main} only after (CMP) is
proved.

Accordingly, the adversarial status is:
\begin{itemize}
\item H1: false as an identification with the fully branch-marked admissible
closure for $N\ge5$; the normalized graph is a nonfinite modification of the
root-stable quotient at the two-scale boundary.
\item H2: partial; the abstract quotient-stack assertion is correct, while
the normalized-Stein, contraction, and conductor comparisons are not proved
globally.
\item H3: plausible but high-risk; the local coarse chart and the finite-map
DVR lemma are correct, but their application to the distinguished repository
component depends on the missing global comparison.
\item OP-LR-NE: open; it requires properness of a specific decorated
$\Isom$ relation, which does not follow from H3 or from ambient properness.
\end{itemize}

This is not a counterexample to the intended theorem.  It is a separation of
the theorem's elementary valuative core from the global moduli comparison
that must be supplied before the theorem can carry downstream dependencies.

\appendix
\section{Explicit pair, triple, and mixed-cluster traits}

This appendix checks the valuative mechanism in coordinates.  These examples
are not substitutes for the global comparison package, but they prevent the
local statement from hiding a residue-field or nilpotent ambiguity.

\subsection{A pair collision}

Let $R=k[[\pi]]$ and consider
\[
 F(T)=(T-a-\pi)(T-a+\pi)
     =T^2-2aT+(a^2-\pi^2).
\]
The unmarked coefficient trait meets the discriminant at $\pi=0$.  The
marked cover is
\[
 R[T]/((T-a)^2-\pi^2).
\]
It has two $R$-sections, $T=a+\pi$ and $T=a-\pi$, corresponding to the two
generic roots.  Both have the same special geometric value $a$.  If the
generic distinguished root is $a+\pi$, separatedness rules out the other
section because they differ over $K$.  Thus special geometric coincidence
does not erase the generic label.

The stack picture contains more information.  At the collision the
unmarked configuration has a transposition exchanging the two colliding
labels.  Choosing one label replaces this stabilizer by the subgroup fixing
that label, which is trivial inside the transposition group.  The stack map
remains etale: the ramification visible in the coarse equation is absorbed
by this change of inertia.  Consequently ``finite etale'' is correct for the
stack map and false for its coarse counterpart.

\subsection{A triple collision}

Let
\[
 F(T)=(T-a-\pi)(T-a)(T-a+\pi).
\]
Writing $U=T-a$ gives $F=U^3-\pi^2U$.  The three generic markings extend as
$U=\pi,0,-\pi$.  The total-collision fiber is $k[U]/(U^3)$, with one
geometric point and length three.  Again there are three different
scheme-valued sections of the family, but no two extend the same generic
section.  At the symmetric special object the unmarked permutation group is
$S_3$, while the stabilizer of a selected label is $S_2$; the relative
inertia kernel remains trivial.

A cyclic transverse presentation such as $U^3=\epsilon$ behaves differently
over the unextended DVR $k[[\epsilon]]$: it has no generic $K$-rational root
in general.  After the ramified base change $\epsilon=\delta^3$ it acquires
the roots $\delta,\zeta\delta,\zeta^2\delta$ when the cube roots of unity are
present.  This does not contradict Lemma~\ref{lem:finite-dvr}.  That lemma
starts with a $K$-rational generic lift.  Properness of a stack may produce a
lift only after a finite base change, whereas H3's strong wording asserts an
extension over the given DVR.  The distinguished normalized root $1$ is
$K$-rational on the clean marked seed, so the stronger conclusion is
available only after the marked finite cover has been globally identified.

\subsection{Two simultaneous clusters}

Let
\[
 F(T)=\prod_{i=1}^{\mu}(T-a-\alpha_i\pi)
       \prod_{j=1}^{\nu}(T-b-\beta_j\pi),\qquad a-b\in R^\times.
\]
Hensel factorization separates the degree-$\mu$ and degree-$\nu$ factors
over $R$.  A generic root in the $a$-cluster cannot specialize into the
$b$-cluster because the two factors define open-and-closed pieces of the
finite root cover.  Inside the $a$-piece all $\mu$ geometric roots can merge,
leaving special algebra $k[U]/(U^\mu)$ after a total collision.  The
distinguished section still exists uniquely because its generic value is
fixed.  This proves the multicluster assertion without pretending that the
entire special root cover has one geometric point: it has at least the two
points $a$ and $b$.

The same argument works when the cluster centers have nontrivial residue
fields.  After a finite separable extension of the residue field, Hensel
factorization produces geometric factors permuted by Galois.  A
$K$-rational selected root defines a Galois-stable section, and the
integral-closure argument descends it to $R$.  If the root is defined only
after a field extension, the conclusion is correspondingly a section only
after normalization in that extension; one must not call it an $R$-valued
mark.

\section{Stack points, coarse points, and sections}

For reference, we spell out the four notions of uniqueness used in the
audit.

\begin{enumerate}
\item \emph{One geometric coarse point} means that the reduced special fiber
of a coarse finite scheme is a singleton.  It ignores nilpotents, tangent
directions, and all generic information.
\item \emph{One geometric stack lift} means one isomorphism class in a
geometric fiber.  Its object can still have a nontrivial automorphism group.
\item \emph{A unique scheme-valued section} means that two morphisms
$\Spec R\to Y$ over the base which agree on $\Spec K$ are equal.  This is a
separatedness statement.
\item \emph{A unique stack-valued lift} means uniqueness up to a specified
2-isomorphism, together with uniqueness of that 2-isomorphism.  The second
part requires trivial relative inertia, which follows here from
representability.
\end{enumerate}

These implications do not reverse in general.  The algebra
$k[T]/(T^\mu)$ has one geometric point but remembers length $\mu$.  The
classifying stack $BG$ has one geometric isomorphism class over an
algebraically closed field but nontrivial automorphisms.  A separated finite
scheme can have many sections with different generic values.  H3 needs the
third notion, while its current last paragraph explicitly proves only the
first and then invokes categorical coarse descent for the rest.

In the quotient model all four levels can be controlled.  The stack map is
representable, so relative automorphisms vanish.  The coarse map is finite,
so a fixed rational generic section extends uniquely.  The invariant ring
computes its nonreduced special fiber.  What remains missing is not a fifth
notion of uniqueness, but the global theorem saying that this quotient model
is the repository's chosen Hurwitz contraction.

\section{Why completion does not supply the missing global morphism}

Suppose $Y$ and $Z$ are finite over a normal base $S$, and suppose their
completed local rings happen to be isomorphic at corresponding closed
points.  To conclude $Y\cong Z$ one normally needs a global $S$-morphism
$Y\to Z$ (or compatible formal isomorphisms satisfying effectivity and
descent), plus hypotheses such as finiteness and birationality.  An
unindexed collection of local isomorphisms does not specify how generic
components are permuted and need not glue on overlaps.

The same warning applies to conductor squares.  For a finite birational
extension $A\subset\ol A$, the conductor is
\[
 \mathfrak c_A=\{x\in\ol A:x\ol A\subset A\}.
\]
It depends on the embedded subring $A\subset\ol A$, not only on $\ol A$ and
the valuation of an ideal on each branch.  Two different subrings of one
normalization may have the same conductor exponent while imposing different
gluing conditions on residue classes.  Therefore the multicluster exponent
calculation in H2 is valuable evidence but cannot identify the full
normalization--conductor pushout until the contracted downstairs algebra has
been identified.

There is a clean route once a morphism exists.  If $Y\to Z$ is finite
birational and $Z$ is normal, then it is an isomorphism.  If both are finite
normal $S$-schemes and a generic isomorphism identifies them with the
integral closure of $S$ in the same fixed generic algebra, that generic
isomorphism extends uniquely.  In both arguments the phrase ``same generic
algebra'' includes an actual identification, and the global morphism or
common embedding is essential.  This is exactly what item 6 of (CMP) asks
H2 to construct.

\begin{thebibliography}{9}
\bibitem{ACV}
D. Abramovich, A. Corti, and A. Vistoli,
\emph{Twisted bundles and admissible covers}, Comm. Algebra 31 (2003),
3547--3618, \href{https://arxiv.org/abs/math/0106211}{arXiv:math/0106211}.

\bibitem{Deopurkar}
A. Deopurkar, \emph{Compactifications of Hurwitz spaces}, Int. Math. Res.
Not. (2014), \href{https://arxiv.org/abs/1206.4535}{arXiv:1206.4535}.

\bibitem{StacksFinite}
The Stacks Project, chapters on finite morphisms, quotient stacks, coarse
moduli spaces, and Stein factorization,
\href{https://stacks.math.columbia.edu}{stacks.math.columbia.edu}.

\bibitem{FP}
B. Fantechi and R. Pandharipande, \emph{Stable maps and branch divisors},
Compos. Math. 130 (2002), 345--364,
\href{https://arxiv.org/abs/math/9905104}{arXiv:math/9905104}.
\end{thebibliography}

\end{document}
